% Preface to the foundations-first second edition.
\chapter*{How to read this book}
\addcontentsline{toc}{chapter}{How to read this book}

You type a sentence. A model returns another token. Between those two
events, software turns text into integers, reads learned numbers,
evaluates a decoder, chooses an output and keeps enough state to do it
again. This book opens those steps one at a time.

The destination is an inference service that is correct, fast and
affordable. We will get there by building understanding before adding
scale. You do not need to know a serving framework, a GPU programming
language or the vocabulary of performance engineering before beginning.
Ordinary programming experience and a willingness to work through
small arrays are enough to start.

\section*{Part I is the main path}
Part~\ref{part:foundations} begins with text and tokens, builds a small
numerical language model, and then constructs the Transformer operators
and a complete decoder. You will calculate attention weights, trace
residual additions and compare cached generation with recomputation.
After that, we examine model files, numerical precision, CPUs, GPUs,
units and reproducible measurements.

These chapters were previously placed in appendices. They now lead the
book because they explain the objects whose cost the later chapters
measure. Their slower, concrete style is the model for the rest of the
manuscript: ask one question, work a small example, draw what changes,
run the code, then test what could go wrong.

Experienced readers can skim the construction and begin with
Part~\ref{part:physics}. The foundation references in later chapters
point back to the relevant explanation; they are not prerequisites
hidden at the back of the book.

\section*{Then make the model an engine}
The remaining parts ask what happens when a correct decoder must
answer quickly, share hardware with other requests and survive failures.
We study memory traffic, prefill and decode, KV storage, optimized
kernels, batching and scheduling. We then examine speculation, changing
model architectures, distributed execution, constrained hardware,
correctness, measurement and production operation.

Four resources organize those engineering choices: memory bandwidth,
memory capacity, computation and latency. A small table records what
each technique saves and what it costs. The table comes after the
mechanism, not in place of it. Correct answers and useful work at an
acceptable cost remain the goals.

\section*{Build, predict and check}
The teaching approach is original, while sharing the build-to-understand
principle of Sebastian Raschka's \emph{Build a Large Language Model
(From Scratch)} \cite{raschka2024build}. The examples here lead to
inference execution and serving, not a training curriculum.

Early chapters use the existing Rust teaching engine and independent
oracles. The new attention, FFN and decoder lessons use standard-library
Python so the arithmetic and state transitions remain visible. Their
weights are illustrative, not trained; passing their tests establishes
computation, not language quality. Later labs return to realistic model
files, kernels and request workloads.

Try to predict each small result before running it. Then change one
assumption: a tensor stride, a future token, a scale, a block size or
a timing endpoint. A useful explanation lets you predict not only
what works but why a nearby version fails.

\section*{How to read the systems drawings}
The drawings answer different questions. A sequence diagram asks who talks
to whom; a state machine asks what may happen next. A component layout
shows where a responsibility or resource lives. A flowchart explains a
decision. We use a small, explained subset of UML notation
\cite{omg2017uml}, alongside tensor and memory drawings:
\begin{center}
\begin{tabularx}{\linewidth}{@{}p{28mm}X@{}}
\toprule
Drawing & Reading rule\\
\midrule
Sequence & Read downward in time. Boxes name participants; vertical dashed
lines are lifelines. Open message arrowheads denote asynchronous signals;
dashed message arrows return a result. Spacing is not elapsed time.\\
State & Rounded boxes are states. A filled circle starts the path and a
bullseye ends it. Brackets mark a guard; text after a slash is an action.\\
Components & Follow the labeled data or control dependency. A box owns
only the responsibility stated in it; it need not be a process or machine.\\
Activity & Follow the steps and the labeled branches of a diamond.
A return arrow means another iteration, not a second independent request.\\
\bottomrule
\end{tabularx}
\end{center}
Captions identify whether a drawing is conceptual, follows the teaching
code or describes a named implementation. An arrow does not prove that
communication is free, asynchronous or already implemented. The nearby
text walks through the decisive path and explains its limits. Sequence views
abstract away the host's blocking calls and thread implementation; logical
message order does not establish CPU/GPU overlap.

\section*{Numbers, identifiers and draft status}
Measured numbers come with setup and raw records. Derived numbers show
their arithmetic. Illustrative numbers are chosen to teach, not to
advertise performance. Current industrial claims cite primary sources
and carry a date; a paper's result is not a measurement made by this book.

Printed chapter numbers follow the new reading order. Code directories
and historical records retain stable identifiers. For example,
\texttt{ch14-request-lifecycle} now accompanies
\Chref{request-lifecycle}; do not infer a printed number from its path.
The source-to-chapter map is in \texttt{docs/STRUCTURE.md}.

The foundational chapters form a complete construction and orientation
path. Many later systems chapters still need their full-depth pass.
The repository status table distinguishes developed chapters from
drafts and independent review from the author's own checks. Read the
book as a working textbook whose claims should remain inspectable
as both the manuscript and the industry evolve.
